% Jacobian Matrix and Velocity Relationship
\documentclass[border=10pt]{standalone}
\usepackage{amsmath}
\pagecolor{black!90}
\color{white}
\begin{document}

\[
J(\mathbf{q}) = 
\begin{bmatrix}
J_v \\ J_\omega
\end{bmatrix}
=
\begin{bmatrix}
\dfrac{\partial \mathbf{p}}{\partial \theta_1} & \cdots & \dfrac{\partial \mathbf{p}}{\partial \theta_6} \\[6pt]
\dfrac{\partial \mathbf{o}}{\partial \theta_1} & \cdots & \dfrac{\partial \mathbf{o}}{\partial \theta_6}
\end{bmatrix}
\qquad (6 \times 6)
\]

\bigskip

\textbf{Velocity mapping:}

\[
\begin{bmatrix} \mathbf{v} \\ \boldsymbol{\omega} \end{bmatrix}
= J(\mathbf{q})\,\dot{\mathbf{q}}
\]

\qquad\qquad

$\mathbf{v}$ = linear velocity, $\boldsymbol{\omega}$ = angular velocity, $\dot{\mathbf{q}}$ = joint velocities

\bigskip

\textbf{Singularity condition:}

\[
\det\!\bigl(J(\mathbf{q})\bigr) = 0
\]

\end{document}